\section{Conclusion}
\label{sec:conclusion}

Medication reconciliation is one of the highest-stakes tasks at clinical handoff points.
Missed or duplicated medications cause preventable patient harm, and every hospital
discharge or transfer is an opportunity for one to slip through. As LLMs begin to be
deployed for this task on FHIR-structured records, the question of how to present FHIR
data to a model has gone largely unexamined. This work provides the first systematic
comparison.

We evaluated 5 models across 4 serialisation strategies on 200 synthetic patients,
totalling 4,000 runs. The findings are consistent and clear. Serialisation strategy
has a statistically significant effect on performance for all instruction-tuned models
up to 8B parameters, with Clinical Narrative being the best-performing format at that
scale. The ranking flips at 70B: Raw JSON becomes the best format, and the strategy
effect loses significance. In all 20 model and strategy combinations, mean precision
exceeds mean recall: omission is the dominant failure mode, and that changes how safety
auditing of these systems should work. Small models plateau at 7--10 active medications;
polypharmacy patients remain systematically underserved below 70B scale. BioMistral's
complete failure adds a fifth finding: domain pretraining without instruction tuning is
not sufficient for structured extraction tasks.

The practical guidance from this work is straightforward. For teams using a 7B--8B
instruction-tuned model, Clinical Narrative serialisation is the right default. For
70B deployments, Raw JSON is simpler and marginally better. In both cases, the key
metric to track in production is recall, not precision. The system's completeness is
what matters clinically, not its accuracy on what it does report.

Future work should extend this benchmark to real hospital EHR data, evaluate
properly instruction-tuned biomedical models, and explore whether few-shot examples
or chain-of-thought prompting can push recall higher for complex patients. The pipeline
uses only open tools: Synthea for patient generation, Ollama for local inference, all
running on an AWS \texttt{g6e.xlarge} instance with an NVIDIA L40S (48 GB VRAM),
making the entire benchmark reproducible without proprietary APIs.
